\clearpage
\section{面試十題：30 秒版本}

每題先講結論，再講一個板上的數字或證據。完整版在括號中的章節。

\begin{qashort}{1. driver 版本比較快，是哪個改動造成的？（第 1 節）}
是 CPU 端的記憶體屬性。driver 用 \code{dma_mmap_coherent} 映射成 Normal non-cacheable，UIO 經 \code{/dev/mem} 是 Device-nGnRnE。四組 A/B 實驗裡，只有最後一步「只改屬性」有差：\abDrvDev{} 到 \abDrvNc{} fps；輪詢位置與 buffer 來源是一起換的，合起來看不出影響。機制是寫入可以合併，CPU 寫 buffer 快 3 到 5 倍，pipeline 每幀「寫輸入」從 0.14--0.15 ms 降到 0.09--0.10 ms。
\end{qashort}

\begin{qashort}{2. 你的 driver 有多少 overhead？（第 2 節）}
空 ioctl 一次 \ioctlUs{} µs；一幀的系統呼叫加參數複製 \bbSysUs{} µs，不到 0.1\%。真正的成本是輪詢：driver 比 PL 多等約 30 到 40 µs，CPU 每幀花 0.45 ms。同一次開機切換比較，改中斷後 CPU 從 \pollCpuUs{} µs 降到 \irqCpuUs{} µs、等待從 \pollWaitUs{} µs 降到 \irqWaitUs{} µs。
\end{qashort}

\begin{qashort}{3. 怎麼把輪詢改成中斷？（第 3 節）}
RTL 加一個暫存器輸出的 level 中斷，等於致能 AND done，致能放在獨立暫存器，免得被 START 寫掉；IP 宣告成 interrupt 介面，接 \code{pl_ps_irq0}，device tree 是 SPI 89。driver 用 \code{devm_request_irq}，ISR 清 done、讀回確保清除到達、\code{complete()}；\code{reinit_completion} 放在 START 之前。1000 幀剛好 1000 次中斷；同一次開機交替量，pipeline 從 \hPollFps{} 變成 \hIrqFps{} fps。
\end{qashort}

\begin{qashort}{4. 改了 RTL 之後 timing 收不了怎麼辦？（第 4 節）}
先固定配方，只動 RTL。建置腳本在 route 後若 WNS 小於 0，自動跑 \code{phys_opt_design}，不行再重新 route。這次把 AxPROT[1] 從 0 改成 1，只改常數，route 完卻是 $-$0.032 ns，靠這一步補到 0。我也準備了備案：Vivado 對同樣的輸入、到 placement 為止給出同樣的結果，打算重現上一版收斂的設計，再用 ECO 只改那幾根常數 pin；最後沒有用到。
\end{qashort}

\begin{qashort}{5. 除錯功能放哪裡？（第 5 節）}
放 debugfs：\code{debugfs_create_regset32} 一次印出暫存器，\code{DEFINE_DEBUGFS_ATTRIBUTE} 做可寫的 \code{nblk}，寫入時檢查範圍、拿 backbone 的鎖。它不是 ABI，app 不該依賴；寫 99 會回 \code{EINVAL}。
\end{qashort}

\begin{qashort}{6. driver 怎麼開機自動載入？（第 6 節）}
Yocto 的 module recipe：\code{inherit module} 用映像自己的 kernel 樹編，vermagic 一定對；\code{KERNEL_MODULE_AUTOLOAD} 讓 systemd 開機載入；device tree 把 \code{compatible} 改成 driver 認得的字串，不再需要 \code{driver_override}。板上開機第 4.8 秒就綁好。
\end{qashort}

\begin{qashort}{7. 為什麼限制 /dev/mem？（第 7 節）}
它把整個實體位址開給 root。driver 用 DMA API 配 buffer、只 mmap buffer，所以 app 不再需要 \code{/dev/mem}，256 MB 保留區也拿掉了。\code{STRICT_DEVMEM} 擋 kernel 管理的 System RAM，\code{IO_STRICT_DEVMEM} 擋 driver 佔用中的暫存器；板上 \code{devmem} 都回 \code{Operation not permitted}。邊界也量過：卸載 driver 後暫存器又讀得到，System RAM 以外的 DDR 最上面 1 MB 也讀得到。
\end{qashort}

\begin{qashort}{8. SMMU 怎麼讓 PL 只能碰自己的 buffer？（第 8 節）}
device tree 打開 SMMU，給加速器的四個 HP 埠登記 stream ID；DMA API 就自動把 buffer 映射到 IOVA，帶著這些 ID 的 Non-secure 交易只能存取映射過的頁，越界讀取會變成 context fault。第一次失敗是因為 AXI master 送的是 Secure 交易，Linux 只設定 Non-secure 那一側；改成 AxPROT[1] = 1 才生效。限制：發佈的 kernel 對沒登記的 stream ID 是 bypass，非 0 的 AXI ID 或 Secure 交易都擋不住。
\end{qashort}

\begin{qashort}{9. 怎麼確定 device tree 節點寫對了？（第 9 節）}
寫 dt-schema 的 binding，規定 v1.1 一定要有中斷、v1.0 不准有。官方檢查要裝 dtschema，這次環境不能裝，所以用自寫的部分檢查比對 driver 的比對表、編譯範例、檢查 PetaLinux 產生的實際節點，因此發現要描述 \code{interrupt-names}。
\end{qashort}

\begin{qashort}{10. 發佈映像怎麼處理密碼？（第 10 節）}
不放真正的密碼：用公開的預設密碼，建置時就用 Yocto 的 \code{passwd-expire} 讓它過期，所以序列埠與 ssh 的第一次登入都要先改；改之前 ssh 帶指令和 scp 都被拒絕，序列埠上 Ctrl-C、空密碼、兩次不一致都拿不到 shell。第一版用 \code{profile.d} 腳本，擋不住非互動的 ssh，才改掉。fixture 換成公式產生的合成圖，所以發佈檔沒有個人資料。限制是 initramfs，重開機就回到過期的預設密碼。
\end{qashort}
